\section{Comandi per l'Architettura di
Rete}\label{comandi-per-larchitettura-di-rete}

Per andare a fondo nei dettagli implementativi e spiegare meglio alcuni
concetti, sono riportati alcuni passaggi fondamentali nella
configurazione dell'host che esegue BlueAgent.

\subsection{configurare OVN}\label{configurare-ovn}

Qui riporto la configurazione per implementare la SDN. Il \textbf{piano
dei dati} e' implementato da Open vSwitch (\textbf{OVS}) e deve esser
configurato in modo che possa comunicare con Incus.

Si configura OVS in modo che possa comunicare con l'interfaccia south
bound di OVN (\passthrough{\lstinline!ovnsb\_db.sock!}):

\begin{lstlisting}[language=bash]
sudo ovs-vsctl set open_vswitch . \ 
external_ids:ovn-remote=unix:/run/ovn/ovnsb_db.sock \ external_ids:ovn-encap-type=geneve \ 
external_ids:ovn-encap-ip=127.0.0.1
\end{lstlisting}

\begin{itemize}
\tightlist
\item
  \textbf{\passthrough{\lstinline!ovn-encap-type=geneve!}}:
  specifichiamo geneve come protocollo da utilizzare per il tunneling.
\item
  \textbf{\passthrough{\lstinline!ovn-encap-ip=127.0.0.1!}}: imposta la
  terminazione del tunnel nel localhost.
\end{itemize}

Si configura invece l'interfaccia north bound in modo che possa
comunicare con Incus:

\begin{lstlisting}[language=bash]
ovn-nbctl set-connection ptcp:6641:127.0.0.1
incus config set network.ovn.northbound_connection tcp:127.0.0.1:6641
\end{lstlisting}

\begin{quote}
\begin{itemize}
\tightlist
\item
  \textbf{OVN Northbound}: e' l'interfaccia dove dichiaro le reti che
  voglio creare e le loro proprietà che devono rispettare.
\item
  \textbf{OVN Southbound DB:}~Memorizza le tabelle di inoltro e la
  logica da applicare sul piano dei dati.
\end{itemize}
\end{quote}

\subsection{\texorpdfstring{\(\text{production-}k\) Routing e
Connettività dei Peer (WireGuard e
nftables)}{\textbackslash text\{production-\}k Routing e Connettività dei Peer (WireGuard e nftables)}}\label{textproduction-k-routing-e-connettivituxe0-dei-peer-wireguard-e-nftables}

L'accesso dei partecipanti ai container rete \(\text{production-}k\)
richiede l'applicazione di regole di inoltro esplicite tra
\passthrough{\lstinline!wg-k!} ed \passthrough{\lstinline!br-team-k!},
da applicare nel kernel Linux mediante
\passthrough{\lstinline!nftables!}.

\subsubsection{\texorpdfstring{Inoltro Kernel e Controllo di Flusso
(\texttt{nftables})}{Inoltro Kernel e Controllo di Flusso (nftables)}}\label{inoltro-kernel-e-controllo-di-flusso-nftables}

Le regole per il traffico vengono definite dentro la tabella
\passthrough{\lstinline!inet blueagent!}:

\begin{lstlisting}[language=bash]
sudo nft add table inet blueagent
sudo nft add set inet blueagent team_links '{ type ifname . ifname; flags interval; }'

sudo nft add element inet blueagent team_links { "wg-$k" . "br-team-$k", "br-team-$k" . "wg-$k" }
\end{lstlisting}

All'interno della catena \passthrough{\lstinline!forward!}, il traffico
viene validato verificando l'appartenenza al set
\passthrough{\lstinline!@team\_links!} e scartando esplicitamente le
comunicazioni incrociate tra team diversi. Vengono inserite
esplicitamente per ogni \passthrough{\lstinline!k!} le coppie di
interfacce per cui autorizzare il traffico.

A priori, si vuole impedire l'inoltro dei pacchetti applicando l'azione
\passthrough{\lstinline!drop!} tra le interfacce
\passthrough{\lstinline!wg!} e \passthrough{\lstinline!br-team!}, in
modo da impedire rotte che non sono state esplicitamente accettate nella
tabella:

\begin{lstlisting}[language=bash]
sudo nft add rule inet blueagent forward iifname "wg-*" oifname "br-team-*" drop
sudo nft add rule inet blueagent forward iifname "br-team-*" oifname "wg-*" drop
sudo nft add rule inet blueagent forward iifname "br-team-*" oifname "br-team-*" drop
sudo nft add rule inet blueagent forward iifname "wg-*" oifname "wg-*" drop
\end{lstlisting}

\subsubsection{Connettività Egress
(Masquerading)}\label{connettivituxe0-egress-masquerading}

Per consentire l'accesso a Internet ai container e ai peer, viene
applicato il mascheramento dell'indirizzo sorgente verso l'esterno
dell'envelope. Si abilita il \passthrough{\lstinline!nat!} (dunque, il
meccanismo di traduzione degli indirizzi) poco prima che il frame esca e
venga immesso fisicamente sulla porta. Poi specifichiamo che si applica
il nat dinamico (\passthrough{\lstinline!masquerade!}) per gli indirizzi
che provengono dalla rete dei container verso l'esterno.

\begin{lstlisting}[language=bash]
sudo nft add chain inet blueagent postrouting '{ type nat hook postrouting priority srcnat; policy accept; }'
sudo nft add rule inet blueagent postrouting ip saddr 10.0.0.0/8 ip daddr != 10.0.0.0/8 masquerade
\end{lstlisting}
